% !TEX program = xelatex
% basic document config
\documentclass[a4paper,10pt,xetex]{article}
\usepackage[a4paper,top=45mm,right=20mm,bottom=30mm,left=25mm,head=35mm,foot=20mm]{geometry}

% language
\usepackage[german]{babel}

% variable definitions
\providecommand{\documenttitle}{Design}
\providecommand{\documentauthors}{Andreas Saurer \\ Benjamin Schneidinger \\ Josef Erben \\ Raffaele Bof \\ Nicolas Loth}
\providecommand{\documentdate}{25.04.2017}

% special commands
\newcommand*{\fullref}[1]{\hyperref[{#1}]{\nameref*{#1} (\ref*{#1})}}
\newcommand{\specialcell}[2][c]{%
  \begin{tabular}[#1]{@{}l@{}}#2\end{tabular}}

% font
% \usepackage{titling}
\usepackage[sfdefault]{roboto}
\usepackage[parfill]{parskip}

% table
\usepackage{array}
\usepackage{multicol}
\usepackage{longtable,tabu,booktabs}

% links
\usepackage{hyperref}
\hypersetup{
            pdftitle={\documenttitle},
            pdfauthor={\documentauthors},
            colorlinks=true,
            linkcolor=[RGB]{74,144,226},
            citecolor=[RGB]{74,144,226},
            urlcolor=[RGB]{74,144,226},
            breaklinks=true}
\urlstyle{same}  % don't use monospace font for urls

% images
\usepackage{float,graphicx,grffile}
\graphicspath{ {images/} }

\makeatletter
\def\maxwidth{\ifdim\Gin@nat@width>\linewidth\linewidth\else\Gin@nat@width\fi}
\def\maxheight{\ifdim\Gin@nat@height>\textheight\textheight\else\Gin@nat@height\fi}
\makeatother
\setkeys{Gin}{width=\maxwidth,height=\maxheight,keepaspectratio}

\makeatletter
\def\fps@figure{H}
\makeatother

% header and footer
\usepackage{lastpage}
\usepackage{fancyhdr}
\pagestyle{fancy}
\fancyhf{}
\fancyhead[L]{\includegraphics[height=2cm]{travel-buddy_white}}
\fancyfoot[L]{\fontsize{8}{10}\selectfont\ \documenttitle}
\fancyfoot[R]{\fontsize{8}{10}\selectfont\ Seite\ \thepage\ von\ \pageref*{LastPage}}

\renewcommand{\headrulewidth}{0pt}
\renewcommand{\footrulewidth}{0pt}

% style titles
\usepackage{titlesec}
\titlespacing*{\section}{0pt}{1em}{0pt}
\titlespacing*{\subsection}{0pt}{1em}{0pt}
\titlespacing*{\subsubsection}{0pt}{1em}{0pt}

% configure title page
\title{
  \includegraphics[width=7cm]{travel-buddy_white}\\[\bigskipamount]
  Anforderungsanalyse\\[\bigskipamount]
}

\author{\documentauthors}
\date{\parbox{\linewidth}{\centering%
  IT15TA ZH \hspace*{3cm} Gruppe 3\endgraf\bigskip
  \documentdate\endgraf
}}


\begin{document}

% title page
\maketitle\newpage

% table of contents
{
\hypersetup{linkcolor=black}
\setcounter{tocdepth}{3}
\tableofcontents
}

\newpage

% version log
\section{Versionenlog}\label{versionenlog}

\tabulinesep=1.2mm

\begin{longtabu} to \textwidth { | l | l | X[l] | l | }
  \hline
  \textbf{Datum} & \textbf{Version} & \textbf{Änderung} & \textbf{Author} \\
  \hline
  \endhead

  0.0.1 & 21.04.2017 & GUI-Design Text hinzugefügt & Andi\\
  \hline

  0.0.0 & 21.04.2017 & Dokument erstellt & Andi\\
  \hline
\end{longtabu}
\newpage

% start content
\section{Projektmanagement}\label{projektmanagement}
\section{Architektur}\label{architektur}
\section{Design-Klassendiagramm}\label{design-klassendiagram}
\section{Klassenverantwortlichkeiten}\label{klassenverantwortlichkeiten}
\section{Interaktionsdiagramme}\label{interaktionsdiagramme}

\section{GUI-Design}\label{gui-design}
Die App startet auf der Anzeige ``TourList'' auf welchem alle verfügbaren Touren angezeigt
werden. Mit einem Klick auf eine Tour wird die Anzeige ``TourDetailView'' geöffnet und eine
Übersicht über die gewählte Tour wird sichtbar. Über den Button ``Tour starten'' gelangt
man weiter zur Ansicht ``RouteView'' worauf jeweils der Weg zum nächsten POI ersichtlich ist.
Abgeschossen wird eine Tour mit der Anzeige ``TourSummary''.

\newpage
\section{Glossar}\label{glossar}
\begin{longtabu} to \textwidth { | l | X[l] | }
\hline
\textbf{Begriff} & \textbf{Bemerkung}\\\hline
\endhead

Agency/Agencies & engl. Agenturen\\\hline
Smartphone & Mobilgerät mit erweiterten Funktionen\\\hline
Android & Betriebssystem für Mobilgeräte\\\hline
Mobile Application & Mobile Applikation, Synonym für mobile Applikation auf Smartphone\\\hline
Backend & Serverseitige Anwendung, stellt Rest API zur Verfügung\\\hline
External Services & engl. Externe Dienste, Dienste die nicht vom Programm selber ausgeführt werden wie z.B. Location Services\\\hline
Fallback & Ausweichlösung/Alternativlösung\\\hline
Frontend & Benutzerseitige Anwendung, hier App welche Rest API konsumiert\\\hline
GPS & Global Positioning System, globales Satelliten Navigationsnetz\\\hline
Marker & Symbol, welches einen Punkt auf der Karte markiert\\\hline
Natives Features & Funktionen, die ein Gerät ohne weitere Software besitzt\\\hline
ÖV & abkz. Öffentlicher Verkehr\\\hline
Point of Interest, PoI & engl. Ort von besonderem Interesse, bspw. Wahrzeichen\\\hline
Stakeholders & engl. Interessenten\\\hline
Tech Stack & Eingesetzte Technologien\\\hline
Mockdaten & Daten mit gleicher Struktur wie echte Daten, dient der lokalen Entwicklung, Testdaten \\\hline
Trip & engl. Reise, enthält mehrere Touren\\\hline
Tourist & Der eigentliche Benutzer der App\\\hline
\end{longtabu}


\section{Literatur}\label{literatur}
\begingroup
\renewcommand{\section}[2]{}%
  \begin{thebibliography}{9}
    \bibitem{UP} C. Larman, Applying UML and patterns. 4. Auflage, Upper Saddle River: Pearson Education, Inc., 2005, S. 18.
  \end{thebibliography}
\endgroup

\section{Abbildungsverzeichnis}\label{abbildungsverzeichnis}
\begingroup
\renewcommand{\section}[2]{}%
\hypersetup{linkcolor=black}
\listoffigures
\endgroup

\end{document}
